% !TEX program = lualatex
\documentclass[12pt,a4paper]{article}

\usepackage{geometry}
\geometry{margin=2cm}
\usepackage{fontspec}

% Imposta HarfBuzz come motore di rendering OpenType moderno in LuaLaTeX
\defaultfontfeatures{Renderer=HarfBuzz}

% Font di base
\setmainfont{Noto Sans}

% Font specifici per scritture del mondo
\newfontfamily\arabicfont[Script=Arabic]{Noto Sans Arabic}
\newfontfamily\hebrewfont[Script=Hebrew]{Noto Sans Hebrew}
\newfontfamily\devanagarifont[Script=Devanagari]{Noto Sans Devanagari}
\newfontfamily\bengalifont[Script=Bengali]{Noto Sans Bengali}
\newfontfamily\thaifont[Script=Thai]{Noto Sans Thai}
\newfontfamily\japanesefont{Noto Sans CJK JP}
\newfontfamily\chinesesimplifiedfont{Noto Sans CJK SC}
\newfontfamily\chinesetraditionalfont{Noto Sans CJK TC}
\newfontfamily\koreanfont{Noto Sans CJK KR}
\newfontfamily\ethiopicfont[Script=Ethiopic]{Noto Sans Ethiopic}
\newfontfamily\georgianfont[Script=Georgian]{Noto Sans Georgian}
\newfontfamily\armenianfont[Script=Armenian]{Noto Sans Armenian}

% Comandi per la direzione nativa Right-to-Left in LuaTeX
\newcommand{\textarabicrtl}[1]{{\textdir TRT\arabicfont #1}}
\newcommand{\texthebrewrtl}[1]{{\textdir TRT\hebrewfont #1}}

\begin{document}

\section*{Test Multilingua Universale per myPdfSearch}

\subsection*{1. Italiano (LTR)}
Questo è un testo di prova in lingua italiana.

\subsection*{2. Spagnolo (LTR)}
Este es un texto de prueba en español: ¿Cómo estás? ¡Muy bien!

\subsection*{3. Giapponese (LTR)}
{\japanesefont こんにちは世界！これは日本語のテキスト抽出テストです。}

\subsection*{4. Cinese Semplificato (LTR)}
{\chinesesimplifiedfont 你好世界！这是一个简体中文的文本提取测试。}

\subsection*{5. Cinese Tradizionale (LTR)}
{\chinesetraditionalfont 你好世界！這是一個繁體中文的文本提取測試。}

\subsection*{6. Coreano (LTR)}
{\koreanfont 안녕하세요 세상! 이것은 한국어 텍스트 추출 테스트입니다.}

\subsection*{7. Arabo (RTL)}
\textarabicrtl{مرحبا بالعالم}! \textarabicrtl{هذا اختبار لاستخراج النص العربي}.

\subsection*{8. Ebraico (RTL)}
\texthebrewrtl{שלום עולם}! \texthebrewrtl{זהו מבחן לחילוץ טקסט בעברית}.

\subsection*{9. Hindi - Devanagari (LTR)}
{\devanagarifont नमस्ते दुनिया}! {\devanagarifont यह हिंदी पाठ का एक परीक्षण है।}

\subsection*{10. Bengalese (LTR)}
{\bengalifont হ্যালো বিশ্ব}! {\bengalifont এটি বাংলা পাঠ্য নিষ্কাশন পরীক্ষা।}

\subsection*{11. Thai (LTR)}
{\thaifont สวัสดีชาวโลก}! {\thaifont นี่คือการทดสอบการดึงข้อความภาษาไทย}

\subsection*{12. Russo - Cirillico (LTR)}
Привет, мир! Это тест извлечения текста на русском языке.

\subsection*{13. Greco Moderno (LTR)}
Γειά σου κόσμε! Αυτή είναι μια δοκιμή εξαγωγής ελληνικού κειμένου.

\subsection*{14. Amarico / Ge'ez - Africa Orientale (LTR)}
{\ethiopicfont ሰላም ልዑል}! {\ethiopicfont ይህ የአማርኛ ጽሑፍ ማውጣት ፈተና ነው።}

\subsection*{15. Georgiano (LTR)}
{\georgianfont გამარჯობა მსოფლიო}! {\georgianfont ეს არის ქართული ტექსტის ამოღების ტესტი.}

\subsection*{16. Armeno (LTR)}
{\armenianfont Բարև աշխարհ}: {\armenianfont Սա հայերեն տեքստի արդյունահանման թեստ է:}

\end{document}
